\documentclass[10pt]{article}
\usepackage[T1]{fontenc}
\usepackage[utf8]{inputenc}
\usepackage[french]{babel}
\usepackage{amsmath,amssymb,mathrsfs,mathpazo}
\usepackage{geometry}
\geometry{a4paper,margin=2cm}
\usepackage{graphicx,array,tabularx,tabularray,colortbl,multirow,multicol,booktabs,enumitem,pifont,ulem,fancybox,diagbox,cancel}
\usepackage{pstricks,pstricks-add,pst-plot,pst-tree,pst-node,pst-text,pst-func,pst-math,pst-bspline,pst-3dplot,pst-eucl,pst-solides3d}
\let\pscalculate\undefined
\usepackage{pst-bezier}
\usepackage{tikz,pgfplots,tkz-tab}
\usetikzlibrary{arrows,arrows.meta,patterns,calc,positioning,trees,decorations.markings,shapes.geometric,shapes.misc,shapes.arrows,intersections,decorations.pathmorphing,decorations.pathreplacing,backgrounds,shadows,babel}
\pgfplotsset{compat=1.18}
\usepgfplotslibrary{fillbetween}
\usepackage[np]{numprint}
\usepackage{listings,makecell,pas-tableur,twoopt}
\let\barre\undefined
\usepackage[pstricks]{tablvar}
\definecolor{darkgreen}{RGB}{0,100,0}
\providecommand{\red}{\color{red}}
\providecommand{\blue}{\color{blue}}
\providecommand{\green}{\color{green}}
\providecommand{\black}{\color{black}}
\providecommand{\white}{\color{white}}
\providecommand{\gray}{\color{gray}}
\providecommand{\cyan}{\color{cyan}}
\providecommand{\magenta}{\color{magenta}}
\providecommand{\yellow}{\color{yellow}}
\providecommand{\violet}{\color{violet}}
\providecommand{\orange}{\color{orange}}
\providecommand{\euro}{\texteuro}
\providecommand{\EUR}{\texteuro}
\providecommand{\ieme}{\textsuperscript{e}}
\providecommand{\degres}{\ensuremath{^\circ}}
\providecommand{\degree}{\ensuremath{^\circ}}
\providecommand{\vv}[1]{\overrightarrow{#1}}
\providecommand{\Coord}[1]{#1}
\providecommand{\decosix}{}
\providecommand{\starredbullet}{\ensuremath{\bullet}}
\pagestyle{empty}
\setlength{\parindent}{0pt}
\usepackage{adjustbox}\DeclareUnicodeCharacter{2236}{:}\DeclareUnicodeCharacter{2190}{\ensuremath{\leftarrow}}
\providecommand{\R}{\mathbb{R}}
\providecommand{\N}{\mathbb{N}}
\providecommand{\Z}{\mathbb{Z}}
\providecommand{\Q}{\mathbb{Q}}
\providecommand{\D}{\mathbb{D}}
\providecommand{\C}{\mathbb{C}}
\providecommand{\e}{{\mathrm{e}}}
\providecommand{\vect}[1]{\overrightarrow{#1}}
\providecommand{\vectt}[1]{\overrightarrow{\mathrm{#1}}}
\providecommand{\vvt}[1]{\overrightarrow{\mathrm{#1}}}
\providecommand{\np}[2][]{{#2~#1}}
\providecommand{\vv}[1]{{\overrightarrow{#1}}}
\renewcommand{\d}{{\mathrm{d}}}
\providecommand{\degres}{{^\circ}}
\providecommand{\degre}{{^\circ}}
\providecommand{\blue}{\color{blue}}
\providecommand{\red}{\color{red}}
\providecommand{\enskip}{\quad}
\providecommand{\llbracket}{\lbrack\!\lbrack}
\providecommand{\rrbracket}{\rbrack\!\rbrack}
\providecommand{\euro}{{\text{€}}}
\providecommand{\no}{{\text{n°}}}
\providecommand{\interval}[2]{{[#1;#2]}}
\providecommand{\intervFO}[2]{{[#1;#2[}}
\providecommand{\dotfill}{\ldots}

\providecommand{\ds}{\displaystyle}
\providecommand{\textding}[1]{#1}
\providecommand{\Oij}{\(\mathrm{O};\vec{\imath},\vec{\jmath}\)}
\providecommand{\Oijk}{\(\mathrm{O};\vec{\imath},\vec{\jmath},\vec{k}\)}
\providecommand{\Ouv}{\(\mathrm{O};\vec{u},\vec{v}\)}
\providecommand{\intervFF}[2]{\ensuremath{[#1;#2]}}
\renewcommand{\euro}{\ifmmode\text{\texteuro}\else\texteuro\fi}
\DeclareUnicodeCharacter{00B0}{\ifmmode^{\circ}\else\textdegree\fi}
\DeclareUnicodeCharacter{20AC}{\ifmmode\text{\texteuro}\else\texteuro\fi}
% Énoncés : Éducation nationale. Transcription et corrigés : APMEP, Denis Vergès et contributeurs. Découpage et classement : CoopMaths, Sébastien Lozano, Éric Elter et contributeurs. Adaptation de la présentation : Maths Explorer.
% CC BY-SA : https://coopmaths.fr/www/about/
% https://coopmaths.fr/alea/static/eam/2025/tex/eam_specifique_2025_06_sujet0v2_1.tex
\begin{document}
\noindent\textit{APMEP / CoopMaths. CC BY-SA. Présentation adaptée par Maths Explorer.}\par\medskip
\textbf{Pour cette première partie, aucune justification n'est demandée et une seule réponse est possible par question. Pour chaque question, reportez son numéro sur votre copie et indiquez votre réponse.}

\subsubsection*{Question 1}
On considère $A = \dfrac{1}{2} - \dfrac{1}{2} \times \dfrac{4}{3}$.

\textbf{a.} $A = 0$ \hfill \textbf{b.} $A = -\dfrac{1}{6}$ \hfill \textbf{c.} $A = \dfrac{2}{3}$ \hfill \textbf{d.} $A = -1$

\subsubsection*{Question 2}
Quatre croissants coûtent 6 euros.

Dix croissants coûtent :

\textbf{a.} 60 euros \hfill \textbf{b.} 8 euros \hfill \textbf{c.} 8,50 euros \hfill \textbf{d.} 15 euros

\subsubsection*{Question 3}
Un prix a doublé. Cela signifie que le prix a augmenté de :

\textbf{a.} 50\% \hfill \textbf{b.} 100\% \hfill \textbf{c.} 150\% \hfill \textbf{d.} 200\%

\subsubsection*{Question 4}
À l'issue d'une augmentation de 10\%, un article coûte 110 euros.

Laquelle des quatre propositions suivantes est vraie ?

\textbf{a.} Le prix de l'article avant l'augmentation était égal à 99 euros. \\
\textbf{b.} Le prix de l'article avant l'augmentation était égal à 120 euros. \\
\textbf{c.} Le prix a augmenté de 10 euros. \\
\textbf{d.} Le prix a augmenté de 11 euros.

\subsubsection*{Question 5}
La masse d'un litre d'huile est égale à 900 grammes.

La masse de 750 millilitres de cette huile est égale à :

\textbf{a.} 750 g \hfill \textbf{b.} 0,675 kg \hfill \textbf{c.} 6,75 kg \hfill \textbf{d.} 67,5 g

\subsubsection*{Question 6}
Dans un repère du plan, on considère les points A$(1~;~100)$ et B$(4~;~106)$. On note $m$ le coefficient directeur de la droite (AB). On peut affirmer que :

\textbf{a.} $m = 2$ \hfill \textbf{b.} $m = 0,5$ \hfill \textbf{c.} $m = -2$ \hfill \textbf{d.} $m = -0,5$

\subsubsection*{Question 7}
Dans un repère du plan, on considère la droite $D$ de coefficient directeur $-0,1$, passant par le point A$(0~;~4)$.

On note B le point de la droite $D$ dont l'abscisse est égale à 1.

L'ordonnée du point B est égale à :

\textbf{a.} 3 \hfill \textbf{b.} 3,9 \hfill \textbf{c.} 4,1 \hfill \textbf{d.} 5

\subsubsection*{Question 8}
La forme développée de $(x - 3)(x + 2)$ est :

\textbf{a.} $x^2 - 5x + 6$ \hfill \textbf{b.} $x^2 - x + 6$ \hfill \textbf{c.} $x^2 - x - 6$ \hfill \textbf{d.} $x^2 - 5x - 6$

\subsubsection*{Question 9}
Le volume $V$ d'un cône de hauteur $h$ et de rayon $r$ est $V = \dfrac{1}{3} \pi r^2 h$.

On cherche à isoler $h$. On a :

\textbf{a.} $h = \dfrac{V}{3 \pi r^2}$ \hfill \textbf{b.} $h = \dfrac{\pi r^2}{3V}$ \hfill \textbf{c.} $h = \dfrac{\sqrt{V}}{\pi r}$ \hfill \textbf{d.} $h = \dfrac{3V}{\pi r^2}$

\subsubsection*{Question 10}
On considère la fonction $f$ définie sur $\mathbb{R}$ par $f(x) = -2x^2 + 3x + 1$.

L'image de $-1$ par la fonction $f$ est égale à :

\textbf{a.} 0 \hfill \textbf{b.} 2 \hfill \textbf{c.} $-2$ \hfill \textbf{d.} $-4$

\subsubsection*{Question 11}
On considère la fonction $f$ définie sur $\mathbb{R}$ par $f(x) = 2x^2 - 5x + 3$.

Un antécédent de 0 par la fonction $f$ est :

\textbf{a.} 1 \hfill \textbf{b.} $-1$ \hfill \textbf{c.} 0 \hfill \textbf{d.} 2

\subsubsection*{Question 12}
On considère les deux séries ci-dessous.

Série A : 9 ; 10 ; 10 ; 11. \hfill Série B : 7 ; 10 ; 10 ; 13.

Laquelle des quatre propositions suivantes est vraie ?

\textbf{a.} La moyenne de la série A est strictement supérieure à la moyenne de la série B. \\
\textbf{b.} La moyenne de la série B est strictement supérieure à la moyenne de la série A. \\
\textbf{c.} L'écart-type de la série A est strictement supérieur à l'écart-type de la série B. \\
\textbf{d.} L'écart-type de la série B est strictement supérieur à l'écart-type de la série A.

\end{document}
